\Needspace{11\baselineskip}

\section{Reading List}

\subsection{Course references (from the course specification)}

\begin{longtable}[]{@{}
  >{\raggedright\arraybackslash}p{(\linewidth - 4\tabcolsep) * \real{0.3333}}
  >{\raggedright\arraybackslash}p{(\linewidth - 4\tabcolsep) * \real{0.3333}}
  >{\raggedright\arraybackslash}p{(\linewidth - 4\tabcolsep) * \real{0.3333}}@{}}
\caption{\centering Course references (from the course specification)}\tabularnewline
\toprule\noalign{}
\begin{minipage}[b]{\linewidth}\raggedright
Reference
\end{minipage} & \begin{minipage}[b]{\linewidth}\raggedright
Role in the course
\end{minipage} & \begin{minipage}[b]{\linewidth}\raggedright
Weeks
\end{minipage} \\
\midrule\noalign{}
\endfirsthead
\toprule\noalign{}
\begin{minipage}[b]{\linewidth}\raggedright
Reference
\end{minipage} & \begin{minipage}[b]{\linewidth}\raggedright
Role in the course
\end{minipage} & \begin{minipage}[b]{\linewidth}\raggedright
Weeks
\end{minipage} \\
\midrule\noalign{}
\endhead
\bottomrule\noalign{}
\endlastfoot
\textbf{Farrell, S., \& Lewandowsky, S. (2018).} \emph{Computational
modeling of cognition and behavior.} Cambridge University Press. &
Primary text for the first half: from verbal to formal models, parameter
estimation, model comparison; also Bayesian estimation, hierarchical
models and model families & 1--10, 12--16 \\
\textbf{Lee, M. D., \& Wagenmakers, E.-J. (2014).} \emph{Bayesian
cognitive modeling: A practical course.} Cambridge University Press. &
Primary text for the second half: Bayesian inference, model comparison
and case studies (SDT, memory retention, GCM) & 2, 4, 7--10, 12--14,
16 \\
\textbf{Kruschke, J. K. (2014).} \emph{Doing Bayesian data analysis: A
tutorial with R, JAGS, and Stan} (2nd ed.). Academic Press. & Gentle,
visual introduction to Bayes' rule, MCMC, diagnostics, hierarchical
models and Bayes factors & 9--13 \\
\textbf{Anderson, B. (2014).} \emph{Computational neuroscience and
cognitive modelling.} SAGE. & Context on computational neuroscience;
neural network and learning models & 1, 15 \\
\textbf{Verguts, T. (2022).} \emph{Introduction to modeling cognitive
processes.} MIT Press. & Accessible modeling introduction; optimization
and gradient descent; delta rule and reinforcement learning;
categorization and decisions & 1, 3, 14--15 \\
\textbf{Ghodousian, A. (1401).} \emph{Mabani-ye ravesh-haye mostaghim
dar behine-sazi gheyr-khatti} (Persian). & Direct-search methods in
nonlinear optimization (theory behind Nelder--Mead) & 3 \\
\textbf{Azimifar, Z., \& Ahmadi, E. (1396).} \emph{Shenasayi-ye amari-ye
olgu: yadgiri-ye nazari-ye moghaddamati}, vol.~1 (Persian). &
Statistical learning view: ML and Bayesian estimation; generalization
and the bias--variance trade-off & 4, 6, 9 \\
\end{longtable}

\subsection{Chapter pointers by week}

\begin{quote}
\textbf{Please verify.} The chapter numbers below were assigned from
memory of the books' tables of contents. Those for Farrell \&
Lewandowsky, Lee \& Wagenmakers and Kruschke are believed to be correct
but should be checked against the editions in the library. For Anderson,
Verguts, Ghodousian and Azimifar \& Ahmadi only topics are given;
chapter numbers are marked \textbf{{[}verify{]}}.
\end{quote}

\begin{longtable}[]{@{}
  >{\raggedright\arraybackslash}p{(\linewidth - 2\tabcolsep) * \real{0.5000}}
  >{\raggedright\arraybackslash}p{(\linewidth - 2\tabcolsep) * \real{0.5000}}@{}}
\caption{\centering Chapter pointers by week}\tabularnewline
\toprule\noalign{}
\begin{minipage}[b]{\linewidth}\raggedright
Week
\end{minipage} & \begin{minipage}[b]{\linewidth}\raggedright
Readings
\end{minipage} \\
\midrule\noalign{}
\endfirsthead
\toprule\noalign{}
\begin{minipage}[b]{\linewidth}\raggedright
Week
\end{minipage} & \begin{minipage}[b]{\linewidth}\raggedright
Readings
\end{minipage} \\
\midrule\noalign{}
\endhead
\bottomrule\noalign{}
\endlastfoot
1 & F\&L ch.~1--2; Verguts, introductory chapter {[}verify{]}; Anderson,
introductory chapter {[}verify{]} \\
2 & F\&L ch.~2; L\&W ch.~11 (SDT, preview), ch.~10 (memory retention) \\
3 & F\&L ch.~3; Ghodousian (direct search methods) {[}verify{]}; Verguts
(optimization/gradient descent) {[}verify{]} \\
4 & F\&L ch.~4; L\&W ch.~3; Azimifar \& Ahmadi (ML estimation)
{[}verify{]} \\
5 & F\&L ch.~5 (and ch.~4 on parameter uncertainty) \\
6 & F\&L ch.~10; Azimifar \& Ahmadi (generalization, bias--variance)
{[}verify{]} \\
7 & F\&L ch.~10; L\&W ch.~7 (introduction) \\
8 & Review: F\&L ch.~2--5, 10; L\&W ch.~11 \\
9 & Kruschke ch.~5--6; L\&W ch.~1, 3; MKL ch.~1; F\&L ch.~6; Azimifar \&
Ahmadi (Bayesian estimation) {[}verify{]} \\
10 & L\&W ch.~7; F\&L ch.~11; Kruschke ch.~10; MKL ch.~2 \\
11 & Kruschke ch.~7; F\&L ch.~7; MKL (MCMC/inference engines)
{[}verify{]} \\
12 & Kruschke ch.~7; MKL ch.~2; L\&W ch.~4 \\
13 & L\&W ch.~11; Kruschke ch.~9; F\&L ch.~9; MKL ch.~2, 10 \\
14 & L\&W ch.~17; F\&L (choice response times) {[}verify chapter{]};
Verguts (categorization, decisions) {[}verify{]} \\
15 & Verguts (delta rule, reinforcement learning) {[}verify{]}; F\&L
(neural networks, neuroscience) {[}verify{]}; Anderson (networks,
learning) {[}verify{]} \\
16 & F\&L (using models in psychology) {[}verify{]}; L\&W Part IV; MKL
ch.~9 \\
\end{longtable}

Abbreviations: F\&L = Farrell \& Lewandowsky; L\&W = Lee \& Wagenmakers;
MKL = Martin, Kumar \& Lao.

\subsection{Optional readings outside the course list}

These are cited in individual weeks for students who want the original
sources. All are well-known publications; please check the details
against the publisher's record before distributing.

\begin{itemize}
\tightlist
\item
  Busemeyer, J. R., \& Wang, Y.-M. (2000). Model comparisons and model
  selections based on generalization criterion methodology.
  \emph{Journal of Mathematical Psychology, 44}, 171--189. (Week 6)
\item
  Ernst, M. O., \& Banks, M. S. (2002). Humans integrate visual and
  haptic information in a statistically optimal fashion. \emph{Nature,
  415}, 429--433. (Week 15)
\item
  Heathcote, A., Brown, S., \& Mewhort, D. J. K. (2000). The power law
  repealed: The case for an exponential law of practice.
  \emph{Psychonomic Bulletin \& Review, 7}, 185--207. (Week 5)
\item
  Marr, D. (1982). \emph{Vision.} W. H. Freeman. (Week 1)
\item
  Medin, D. L., \& Schaffer, M. M. (1978). Context theory of
  classification learning. \emph{Psychological Review, 85}, 207--238.
  (Week 14)
\item
  Roberts, S., \& Pashler, H. (2000). How persuasive is a good fit? A
  comment on theory testing. \emph{Psychological Review, 107}, 358--367.
  (Week 6)
\item
  Vehtari, A., Gelman, A., Simpson, D., Carpenter, B., \& Bürkner, P.-C.
  (2021). Rank-normalization, folding, and localization: An improved
  \(\hat R\) for assessing convergence of MCMC. \emph{Bayesian Analysis,
  16}, 667--718. (Week 12)
\item
  Wagenmakers, E.-J., van der Maas, H. L. J., \& Grasman, R. P. P. P.
  (2007). An EZ-diffusion model for response time and accuracy.
  \emph{Psychonomic Bulletin \& Review, 14}, 3--22. (Week 14)
\item
  Wilson, R. C., \& Collins, A. G. E. (2019). Ten simple rules for the
  computational modeling of behavioral data. \emph{eLife, 8}, e49547.
  (Weeks 7, 15, 16)
\end{itemize}
